\subsection{CRITERIA FOR RATIONALITY}

For every subfield L of K, let \(\operatorname{End}_{\mathrm{L}}(\mathrm{K})\) denote the endomorphism ring of K considered as a left vector space over L; if L contains \(K'\) , \(\operatorname{End}_{\mathrm{L}}(\mathrm{K})\) is a subring of \(\operatorname{End}_{\mathrm{K}'}(\mathrm{K})\) . For every subset M of \(\operatorname{End}_{\mathrm{K}'}(\mathrm{K})\) , there exists a largest subfield L of K containing \(K'\) and such that M is contained in \(\operatorname{End}_{\mathrm{L}}(\mathrm{K})\) , namely the set of \(\xi \in K\) such that \(\phi(\xi\eta) = \xi.\phi(\eta)\) for all \(\eta \in K\) and all \(\phi \in M\) (it is immediately verified that this set is a ring and, on the other hand, replacing \(\eta\) by \(\xi^{-1}\eta\) in the preceding relation, we obtain \(\phi(\xi^{-1}\eta) = \xi^{-1}.\phi(\eta)\) when \(\xi \neq 0\) ). We shall call this field the centralizer of M in K and denote it by \(\chi(\mathcal{M})\) . Now let V be a right vector K-space with a \(K'\) -structure \(V'\) . For all \(\phi \in \operatorname{End}_{\mathrm{K}'}(\mathrm{K})\) , there exists one and only one endomorphism \(\phi_{V}\) of the Z-module V such that \(\phi_{\mathbf{V}}(x'.\xi) = x'.\phi(\xi)\) for \(x' \in \mathbf{V}'\) and \(\xi \in \mathbf{K}\) : for, in no. 1, a Z-isomorphism \(\lambda\) of \(\mathbf{V}' \otimes_{\mathbb{K}'} \mathbf{K}\) onto V was defined which maps \(x' \otimes \xi\) to \(x'\) . \(\xi\) and \(\phi_{\mathbf{V}}\) is necessarily equal to \(\lambda \circ (1_{\mathbf{V}'} \otimes \phi) \circ \lambda^{-1}\) .

THEOREM 1. Let \(\mathcal{M}\) be a subset of \(\operatorname{End}_{\mathbb{K}'}(\mathbf{K})\) and \(\mathbf{L} = \chi(\mathcal{M})\) the subfield of \(\mathbf{K}\) the centralizer of \(\mathcal{M}\) .

\begin{enumerate}
\def\labelenumi{(\roman{enumi})}
\item
  Let V be a right vector K-space with a K'-structure. For a vector \(x \in V\) to be rational over L, it is necessary and sufficient that \(\phi_{\mathbf{v}}(x.\eta) = x.\phi(\eta)\) for all \(\phi \in M\) and all \(\eta \in K\) . For a vector sub-K-space W of V to be rational over L, it is necessary and sufficient that \(\phi_{\mathbf{v}}(\mathbf{W}) \subset \mathbf{W}\) for all \(\phi \in M\) .
\item
  Let \(V_{1}\) , \(V_{2}\) be two right vector K-spaces each with a \(K^{\prime}\) -structure. For a K-linear mapping f of \(V_{1}\) into \(V_{2}\) to be rational over L, it is necessary and sufficient that \(f(\phi_{\mathrm{V}_{1}}(x_{1})) = \phi_{\mathrm{V}_{2}}(f(x_{1}))\) for all \(x_{1} \in V_{1}\) and all \(\phi \in M\) .
\end{enumerate}

We first prove assertion (i) for \(x\) . Let \(B\) be a basis of \(V\) rational over \(K'\) and write \(x = \sum_{b \in B} b \cdot \xi_b\) ; then, for \(\phi \in \mathcal{M}\) and \(\eta \in K\) ,

\[
\phi_ {V} (x. \eta) - x. \phi (\eta) = \sum_ {b \in B} b. (\phi (\xi_ {b} \eta) - \xi_ {b}. \phi (\eta))
\]

and therefore, the relations

\[
" \text { for   all } \phi \in \mathcal {M} \text { and   all } \eta \in K, \phi_ {V} (x. \eta) = x. \phi (\eta)"
\]

and

\[
" \text { for   all } \phi \in \mathcal {M}, \text { all } b \in \mathrm{B} \text { and   all } \eta \in \mathrm{K}, \phi (\xi_ {b} \eta) = \xi_ {b}. \phi (\eta)"
\]

are equivalent. The second of these relations means that for all \(b \in \mathbf{B}\) , \(\xi_b \in \chi(\mathcal{M})\) , which proves the first assertion of (i).

We next prove (ii). For \(f\) to be rational over \(L\) , it is necessary and sufficient that, for all \(x_1' \in V_1\) rational over \(K', f(x_1')\) is a vector in \(V_2\) rational over \(L\) ; this will imply that \(f(x_1)\) is rational over \(L\) for every vector \(x_1\) of \(V_1\) rational over \(L\) , such a vector being a linear combination with coefficients in \(L\) of vectors rational over \(K'\) . The above condition is equivalent, by the first part of the argument, to the relation

\[
f (x _ {1} ^ {\prime}) \cdot \phi (\eta) = \phi_ {\mathrm{v} _ {2}} (f (x _ {1} ^ {\prime}) \cdot \eta) \quad \text { for } \phi \in \mathcal {M} \text { and } \eta \in \mathrm{K}\tag{4}
\]

which may also be written

\[
f (\phi_ {\mathrm{v} _ {1}} (x _ {1} ^ {\prime}. \eta)) = \phi_ {\mathrm{v} _ {2}} (f (x _ {1} ^ {\prime}. \eta)) \quad \text { for } \phi \in \mathcal {M} \text { and } \eta \in \mathrm{K}.\tag{5}
\]

As every element of \(V_{1}\) is a linear combination with coefficients in K of elements of \(V_{1}\) rational over \(K'\) , condition (5) is equivalent to

\[
f (\phi_ {\mathrm{v} _ {1}} (x _ {1})) = \phi_ {\mathrm{v} _ {2}} (f (x _ {1}))
\]

for all \(x_{1} \in V_{1}\) and all \(\phi \in \mathcal{M}\) .

Finally, to prove the second assertion in (i), we use no. 6, Lemma 1: W is the graph of a K-linear mapping \(g \colon \mathbb{W}_1 \to \mathbb{W}_2\) and W is rational over L if and only if the mapping \(g\) is rational over L (no. 3, Proposition 4). By (ii), for \(g\) to be rational over L, it is necessary and sufficient that \(g(\phi_{\mathbb{W}_1}(x_1)) = \phi_{\mathbb{W}_2}(g(x_1))\) for all \(x_1 \in \mathbb{W}_1\) and all \(\phi \in \mathcal{M}\) ; as \(\phi_{\mathbb{V}} = \phi_{\mathbb{W}_1} \times \phi_{\mathbb{W}_2}\) , the above condition means that the graph W of \(g\) is stable under \(\phi_{\mathbb{V}}\) for all \(\phi \in \mathcal{M}\) .

